Fix \(m\geq1\), \(p_*\in[\epsilon,1-\epsilon]\), and \(q\in K_{\epsilon}\).  Define
\[
 r=p(q),\qquad y_1=\frac{q_{11}}r,\qquad y_0=\frac{q_{01}}{1-r},
 \qquad w_1=w_m(r,p_*),\qquad w_0=w_m(1-r,1-p_*).
\]
The restrictions give \(r,1-r,p_*,1-p_*>0\), and the simplex inequalities give \(0\leq y_0,y_1\leq1\).  Put \(a_1=1-w_1\) and \(a_0=1-w_0\).  From the definition of \(w_m\),
\[
 a_1=
 \begin{cases}
 (1-r)((p_*-r)/p_*)^{m-1},&m\text{ odd},\\
 ((p_*-r)/p_*)^m,&m\text{ even},
 \end{cases}
\]
and \(a_0\) has the same form after replacing \((r,p_*)\) by \((1-r,1-p_*)\).  In either parity the exponent is even and the remaining multiplier is nonnegative, so
\[
 a_0\geq0,\qquad a_1\geq0.
\tag{1}
\]
If \(m\geq2\) and \(r\neq p_*\), all powered ratios are nonzero and the odd-case multipliers are positive; therefore
\[
 a_0>0,\qquad a_1>0.
\tag{2}
\]

Substitution in the two envelope definitions gives
\[
 \ell_{m,p_*}(q)=y_1w_1-y_0w_0-a_0,
 \qquad
 u_{m,p_*}(q)=y_1w_1-y_0w_0+a_1.
\]
Since \(f(q)=y_1-y_0\),
\[
 \begin{aligned}
 f(q)-\ell_{m,p_*}(q)&=y_1a_1+(1-y_0)a_0,\\
 u_{m,p_*}(q)-f(q)&=(1-y_1)a_1+y_0a_0.
 \end{aligned}
\tag{3}
\]
Every factor on the right is nonnegative by (1), so
\[
 \ell_{m,p_*}(q)\leq f(q)\leq u_{m,p_*}(q)
 \qquad(q\in K_{\epsilon}).
\tag{4}
\]

For fixed \(p_*\), \(r\mapsto w_m(r,p_*)\) is a polynomial of degree at most \(m\) and vanishes at \(r=0\), so it is divisible by \(r\), with quotient of degree at most \(m-1\).  Likewise \(w_m(1-r,1-p_*)\) is divisible by \(1-r\).  Thus the apparent propensity denominators in \(\ell_{m,p_*}\) and \(u_{m,p_*}\) cancel, and both are polynomials of total degree at most \(m\) in the cell coordinates.  Homogenizing a monomial of degree \(d<m\) by multiplication with
\[
 (q_{00}+q_{01}+q_{10}+q_{11})^{m-d}=1
\]
on the simplex expresses it as a linear combination of the degree-\(m\) Bernstein coordinates.  Hence \(\ell_{m,p_*},u_{m,p_*}\in V_m\).  Together with (4), this proves
\[
 (\ell_{m,p_*},u_{m,p_*})\in\mathcal E_m.
\tag{5}
\]

Fix \(b\in\mathcal B_{m,\epsilon}\) and choose \(\nu\in\mathcal F_{m,\epsilon}(b)\).  Since both envelopes lie in \(V_m\), their Bernstein coefficient vectors and the moment equations give
\[
 A(p_*;b):=\ell_{m,p_*}^{\top}b=\int\ell_{m,p_*}\,d\nu,
 \qquad
 H(p_*;b):=u_{m,p_*}^{\top}b=\int u_{m,p_*}\,d\nu.
\tag{6}
\]
These values are independent of the chosen representation.  The displayed formulas are jointly continuous on the compact set \(K_{\epsilon}\times[\epsilon,1-\epsilon]\), because all denominators are uniformly bounded away from zero.  Uniform continuity and (6) show that \(A(\cdot;b)\) and \(H(\cdot;b)\) are continuous.  Consequently
\[
 \max_{p_*\in[\epsilon,1-\epsilon]}A(p_*;b),
 \qquad
 \min_{p_*\in[\epsilon,1-\epsilon]}H(p_*;b)
\tag{7}
\]
are attained.

For each \(p_*\), (5) and the sharp-envelope domination theorem give
\[
 A(p_*;b)\leq L_m(b)\leq U_m(b)\leq H(p_*;b).
\tag{8}
\]
Taking the maximum on the left and minimum on the right proves
\[
 [L_m(b),U_m(b)]\subseteq I_{LW,m}^{\cap}(b).
\tag{9}
\]
Every interval \([A(p_*;b),H(p_*;b)]\) contains the common nonempty interval \([L_m(b),U_m(b)]\).  The endpoint formula for an intersection of intervals and (7) therefore give
\[
 I_{LW,m}^{\cap}(b)
 =\bigcap_{p_*\in[\epsilon,1-\epsilon]}I_{LW,m,p_*}(b).
\tag{10}
\]

Assume \(m\geq2\).  For \(s\in[\epsilon,1-\epsilon]\), define
\[
 \mathcal C_s^-=\{q\in K_{\epsilon}:\ell_{m,s}(q)=f(q)\},
 \qquad
 \mathcal C_s^+=\{q\in K_{\epsilon}:u_{m,s}(q)=f(q)\}.
\]
If \(p(q)\neq s\), (2) makes both \(a_0\) and \(a_1\) positive.  The first identity in (3) can then vanish only if \(y_1=0\) and \(y_0=1\), and the second only if \(y_1=1\) and \(y_0=0\).  Since both propensity denominators are positive,
\[
 \begin{aligned}
 \mathcal C_s^-&\subseteq\{q:p(q)=s\}\cup\{q:q_{11}=q_{00}=0\},\\
 \mathcal C_s^+&\subseteq\{q:p(q)=s\}\cup\{q:q_{10}=q_{01}=0\}.
 \end{aligned}
\tag{11}
\]

For every polynomial \(h\) of degree at most two and every feasible \(b\), define
\[
 M_h(b)=\int h(q)\,d\nu(q),\qquad \nu\in\mathcal F_{m,\epsilon}(b).
\]
This is well defined because \(m\geq2\) and homogenization puts \(h\in V_m\), so its integral is a fixed linear combination of \(b\).  Define
\[
 \begin{aligned}
 D_-(b)&=M_{p q_{11}}(b)M_{q_{00}}(b)-M_{p q_{00}}(b)M_{q_{11}}(b),\\
 D_+(b)&=M_{p q_{10}}(b)M_{q_{01}}(b)-M_{p q_{01}}(b)M_{q_{10}}(b).
 \end{aligned}
\tag{12}
\]
Each \(M_h\) is linear in \(b\), so \(D_-\) and \(D_+\) are polynomials on \(\operatorname{aff}(\mathcal B_{m,\epsilon})\).

Suppose \(L_m(b)=\max_sA(s;b)\), and let \(s_b\) attain the maximum.  The lower equality characterization in the sharp-envelope theorem gives a representing measure supported on \(\mathcal C_{s_b}^-\).  By (11), every support point satisfies \((p-s_b)q_{11}=0\) and \((p-s_b)q_{00}=0\).  Integration gives
\[
 M_{p q_{11}}(b)=s_bM_{q_{11}}(b),
 \qquad M_{p q_{00}}(b)=s_bM_{q_{00}}(b),
\]
and hence \(D_-(b)=0\).  Thus
\[
 L_m(b)=\max_sA(s;b)\Longrightarrow D_-(b)=0.
\tag{13}
\]
The upper equality characterization and the second containment in (11) similarly give
\[
 U_m(b)=\min_sH(s;b)\Longrightarrow D_+(b)=0.
\tag{14}
\]

Put
\[
 q^{(0)}=(1-\epsilon,0,\epsilon,0),
 \qquad q^{(1)}=(0,\epsilon,0,1-\epsilon),
 \qquad \nu^{\dagger}=\tfrac12\delta_{q^{(0)}}+\tfrac12\delta_{q^{(1)}},
\]
and \(b^{\dagger}=\int g_m\,d\nu^{\dagger}\).  Both support points lie in \(K_{\epsilon}\).  Direct substitution in (12) yields
\[
 D_-(b^{\dagger})=\tfrac14(1-\epsilon)^2(1-2\epsilon)>0,
 \qquad
 D_+(b^{\dagger})=\tfrac14\epsilon^2(2\epsilon-1)<0,
\tag{15}
\]
using \(0<\epsilon<1/2\).

Let \(\mathcal A=\operatorname{aff}(\mathcal B_{m,\epsilon})\) and
\[
 \mathcal R=\operatorname{relint}(\mathcal B_{m,\epsilon})
 \cap\{b\in\mathcal A:D_-(b)D_+(b)\neq0\}.
\tag{16}
\]
This is relatively open.  Choose \(b^{\circ}\in\operatorname{relint}(\mathcal B_{m,\epsilon})\) and put \(b_t=(1-t)b^{\dagger}+tb^{\circ}\), \(0<t\leq1\).  If the relative ball of radius \(\rho\) about \(b^{\circ}\) lies in the body, convexity puts the relative ball of radius \(t\rho\) about \(b_t\) in the body.  Thus \(b_t\) is in the relative interior.  By (15), continuity of \(D_-D_+\), and \(b_t\to b^{\dagger}\), the product is nonzero for all sufficiently small \(t>0\).  Hence \(\mathcal R\neq\varnothing\).

For \(b\in\mathcal R\), (13) rules out lower equality in (8), and (14) rules out upper equality.  Therefore
\[
 L_m(b)>\max_{p_*\in[\epsilon,1-\epsilon]}\ell_{m,p_*}^{\top}b,
 \qquad
 U_m(b)<\min_{p_*\in[\epsilon,1-\epsilon]}u_{m,p_*}^{\top}b.
\]
This proves the common nonempty relatively open strict-improvement region.